\paragrafos{Γωνίες σε παράλληλες ευθείες με τέμνουσα}{Β.3.8, «Γωνίες σε παράλληλες ευθείες με τέμνουσα»}

\begin{anagke}
Η θεία Μαρία φυτεύει ντομάτες σε δύο ευθείες σειρές, που είναι παράλληλες, και ένα λοξό μονοπάτι τις διασχίζει και τις δύο. Η Ρήνη, που βοηθά στον κήπο, μετρά με τον μοιρογνωμόνιο τη γωνία ανάμεσα στο μονοπάτι και στην πρώτη σειρά: $62^\circ$. «Να μετρήσω και την άλλη σειρά;», ρωτά. «Δεν χρειάζεται», λέει η κυρία Αλεξάνδρα, που περνούσε τυχαία, «αν ξέρεις τι σημαίνει να είναι παράλληλες, τη βρίσκεις με το μυαλό». Ο Ηλίας μετρά, για να βεβαιωθεί: $62^\circ$. Και στις υπόλοιπες θέσεις; Βγαίνουν συνολικά οκτώ γωνίες και, όπως φαίνεται, μόνο δύο διαφορετικά μέτρα.

Για να μιλήσουμε με ακρίβεια για το ποια γωνία είναι ίση με ποια, χρειαζόμαστε ονόματα: «εντός», «εκτός», «εναλλάξ», «επί τα αυτά».

Σε αυτή την παράγραφο μαθαίνουμε να ονομάζουμε τις οκτώ γωνίες που σχηματίζονται όταν δύο παράλληλες ευθείες τέμνονται από μια τρίτη, και να βρίσκουμε όλες τις γωνίες από μία μόνο.
\end{anagke}

\begin{lexi}{τέμνουσα, εντός, εκτός, εναλλάξ, επί τα αυτά}{θέσεις γωνιών}
Η \lexiword{τέμνουσα} είναι η ευθεία που «κόβει» τις άλλες δύο. Το \lexiword{εντός} σημαίνει «μέσα» και το \lexiword{εκτός} «έξω». Το \lexiword{εναλλάξ} είναι οι θέσεις που «αλλάζουν» πλευρά, δηλαδή αντίθετες. Το \lexiword{επί τα αυτά} σημαίνει «προς τα ίδια μέρη», δηλαδή προς την ίδια πλευρά.
\end{lexi}

\begin{metafrasi}
Η ευθεία που κόβει δύο άλλες ευθείες. & τέμνουσα\\
Μέσα στη λωρίδα ανάμεσα στις παράλληλες. & εντός\\
Έξω από τη λωρίδα των παραλλήλων. & εκτός\\
Στην ίδια πλευρά της τέμνουσας. & επί τα αυτά\\
Σε αντίθετες πλευρές της τέμνουσας. & εναλλάξ\\
\end{metafrasi}

\tmima{Η θεωρία}

Δύο ευθείες $\varepsilon_1$ και $\varepsilon_2$ του ίδιου επιπέδου είναι \textbf{παράλληλες} ($\varepsilon_1\parallel\varepsilon_2$) όταν δεν έχουν κανένα κοινό σημείο, όσο κι αν προεκταθούν. Αν μια τρίτη ευθεία $\varepsilon_3$ τις τέμνει και τις δύο, λέγεται \textbf{τέμνουσά} τους. Τότε σχηματίζονται οκτώ γωνίες, τις οποίες ονομάζουμε $\alpha,\beta,\gamma,\delta$ γύρω από το πρώτο σημείο τομής και $\varepsilon,\zeta,\eta,\theta$ γύρω από το δεύτερο.

\begin{center}
\begin{tikzpicture}[scale=1.05]
\draw[thick] (-3,1.2)--(3,1.2) node[right] {$\varepsilon_1$};
\draw[thick] (-3,-1.2)--(3,-1.2) node[right] {$\varepsilon_2$};
\draw[thick,gblue] (1.4,2.8)--(-1.4,-2.8) node[below left] {$\varepsilon_3$};
\fill (0.6,1.2) circle (1.6pt); \fill (-0.6,-1.2) circle (1.6pt);
\node at (1.05,1.47) {\small $\alpha$}; \node at (0.33,1.65) {\small $\beta$};
\node at (0.15,0.75) {\small $\gamma$}; \node at (1.0,0.75) {\small $\delta$};
\node at (-0.15,-0.93) {\small $\varepsilon$}; \node at (-0.87,-0.75) {\small $\zeta$};
\node at (-1.05,-1.47) {\small $\eta$}; \node at (-0.15,-1.65) {\small $\theta$};
\end{tikzpicture}
\end{center}

\textbf{Ως προς τις παράλληλες:} οι γωνίες $\gamma,\delta,\varepsilon,\zeta$ βρίσκονται ανάμεσα στις $\varepsilon_1,\varepsilon_2$ και λέγονται \textbf{εντός}. Οι $\alpha,\beta,\eta,\theta$ βρίσκονται έξω από τη λωρίδα και λέγονται \textbf{εκτός}.

\textbf{Ως προς την τέμνουσα:} δύο γωνίες που βρίσκονται προς το ίδιο μέρος της $\varepsilon_3$ λέγονται \textbf{επί τα αυτά}. Αν βρίσκονται σε διαφορετικά μέρη, λέγονται \textbf{εναλλάξ}.

Συνδυάζοντας τα δύο, παίρνουμε τα ζεύγη του πίνακα.

\begin{center}
\renewcommand{\arraystretch}{1.25}
\begin{tabular}{@{}p{0.37\linewidth}p{0.29\linewidth}p{0.28\linewidth}@{}}
\toprule
Ζεύγος & Ονομασία & Σχέση\\
\midrule
$\gamma,\varepsilon$ και $\delta,\zeta$ & εντός εναλλάξ & ίσες\\
$\alpha,\eta$ και $\beta,\theta$ & εκτός εναλλάξ & ίσες\\
$\alpha,\varepsilon$ και $\beta,\zeta$ και $\gamma,\eta$ και $\delta,\theta$ & εντός--εκτός και επί τα αυτά & ίσες\\
$\delta,\varepsilon$ και $\gamma,\zeta$ & εντός και επί τα αυτά & παραπληρωματικές\\
$\alpha,\theta$ και $\beta,\eta$ & εκτός και επί τα αυτά & παραπληρωματικές\\
\bottomrule
\end{tabular}
\end{center}

\begin{kanonas}{παράλληλες ευθείες και τέμνουσα}
Αν $\varepsilon_1\parallel\varepsilon_2$, τότε:
\begin{itemize}[leftmargin=1.2em,itemsep=0.3ex,topsep=2pt]
\item οι \textbf{εντός εναλλάξ} γωνίες είναι ίσες,
\item οι \textbf{εκτός εναλλάξ} γωνίες είναι ίσες,
\item οι \textbf{εντός και εκτός επί τα αυτά} γωνίες είναι ίσες,
\item οι \textbf{εντός και επί τα αυτά} γωνίες έχουν άθροισμα $180^\circ$ (είναι παραπληρωματικές), και το ίδιο οι \textbf{εκτός και επί τα αυτά}.
\end{itemize}
Με απλά λόγια: όλες οι \textbf{οξείες} γωνίες του σχήματος είναι ίσες μεταξύ τους, όλες οι \textbf{αμβλείες} είναι ίσες μεταξύ τους, και μια οξεία με μια αμβλεία έχουν άθροισμα $180^\circ$.
\end{kanonas}

Γιατί συμβαίνει αυτό; Όταν δύο ευθείες είναι παράλληλες, η τέμνουσα «γέρνει» προς αυτές με την ίδια κλίση και στις δύο, άρα οι αντίστοιχες γωνίες είναι ίδιες. Οι υπόλοιπες προκύπτουν από τις κατακορυφήν γωνίες (ίσες) και τις παραπληρωματικές (άθροισμα $180^\circ$).

Ισχύει και το \textbf{αντίστροφο}, που μας βοηθά να ελέγχουμε αν δύο ευθείες είναι παράλληλες: αν δύο ευθείες τεμνόμενες από τρίτη σχηματίζουν \emph{εντός εναλλάξ} γωνίες ίσες (ή εντός και επί τα αυτά γωνίες με άθροισμα $180^\circ$), τότε οι ευθείες είναι παράλληλες.

\begin{prosoxi}
\begin{itemize}[leftmargin=1.2em,itemsep=0.3ex,topsep=2pt]
\item Οι ισότητες ισχύουν \emph{μόνο όταν οι ευθείες είναι παράλληλες}. Αν δεν είναι, οι εντός εναλλάξ γωνίες δεν είναι ίσες.
\item Οι εντός και επί τα αυτά γωνίες είναι \emph{παραπληρωματικές}, όχι ίσες (εκτός αν είναι και οι δύο $90^\circ$).
\item Μην μπερδεύεις το «εναλλάξ» με το «επί τα αυτά»: το πρώτο σημαίνει διαφορετικές πλευρές της τέμνουσας, το δεύτερο την ίδια.
\end{itemize}
\end{prosoxi}

\tmima{Λυμένα παραδείγματα}

\paradeigma{Όλες οι γωνίες από μία}
\tixero $\varepsilon_1\parallel\varepsilon_2$ και στο σχήμα της θεωρίας $\gamma=52^\circ$.\par
\tithelo Οι υπόλοιπες επτά γωνίες.\par
\possindeo Η $\alpha$ είναι κατακορυφήν με τη $\gamma$, άρα $\alpha=52^\circ$. Η $\beta$ είναι παραπληρωματική της $\gamma$ (σχηματίζουν ευθεία), άρα $\beta=180^\circ-52^\circ=128^\circ$, και η $\delta=\beta=128^\circ$ (κατακορυφήν). Η $\varepsilon$ είναι εντός εναλλάξ με τη $\gamma$, άρα $\varepsilon=52^\circ$. Η $\zeta$ είναι εντός και επί τα αυτά με τη $\gamma$, άρα $\zeta=180^\circ-52^\circ=128^\circ$. Τέλος $\eta=\varepsilon=52^\circ$ και $\theta=\zeta=128^\circ$ (κατακορυφήν).\par
\apantisi $\alpha=\gamma=\varepsilon=\eta=52^\circ$ και $\beta=\delta=\zeta=\theta=128^\circ$.\par
\elegxos Τέσσερις οξείες ίσες, τέσσερις αμβλείες ίσες, και $52^\circ+128^\circ=180^\circ$.

\paradeigma{Εντός εναλλάξ με άγνωστο}
\tixero $\varepsilon_1\parallel\varepsilon_2$ και δύο εντός εναλλάξ γωνίες είναι $x+20^\circ$ και $2x-40^\circ$.\par
\tithelo Το $x$ και οι δύο γωνίες.\par
\possindeo Οι εντός εναλλάξ γωνίες είναι ίσες, άρα $x+20^\circ=2x-40^\circ$. Προσθέτω $40^\circ$ και στα δύο μέλη: $x+60^\circ=2x$, άρα $x=60^\circ$.\par
\apantisi $x=60^\circ$ και κάθε γωνία είναι $80^\circ$.\par
\elegxos $60^\circ+20^\circ=80^\circ$ και $2\cdot60^\circ-40^\circ=80^\circ$.

\paradeigma{Εντός και επί τα αυτά με άγνωστο}
\tixero $\varepsilon_1\parallel\varepsilon_2$ και δύο εντός και επί τα αυτά γωνίες είναι $4x$ και $x+20^\circ$.\par
\tithelo Οι δύο γωνίες.\par
\possindeo Είναι παραπληρωματικές: $4x+x+20^\circ=180^\circ$, άρα $5x=160^\circ$ και $x=32^\circ$.\par
\apantisi $4x=128^\circ$ και $x+20^\circ=52^\circ$.\par
\elegxos $128^\circ+52^\circ=180^\circ$.

\paradeigma{Παράλληλες ή όχι;}
\tixero Δύο ευθείες τέμνονται από τρίτη και οι εντός εναλλάξ γωνίες είναι $72^\circ$ και $75^\circ$.\par
\tithelo Είναι παράλληλες;\par
\possindeo Αν ήταν παράλληλες, οι εντός εναλλάξ γωνίες θα ήταν ίσες.\par
\apantisi Όχι, γιατί $72^\circ\neq75^\circ$.\par
\elegxos Αν και οι δύο ήταν $72^\circ$, θα ήταν παράλληλες.

\tmima{Ασκήσεις}

\zoni{Το κατάλαβα}{ονόματα και ισότητες}

\ex Στο σχήμα της θεωρίας, πες πώς ονομάζεται κάθε ζεύγος γωνιών: α) $\alpha,\varepsilon$ β) $\gamma,\varepsilon$ γ) $\beta,\theta$ δ) $\delta,\varepsilon$ ε) $\gamma,\zeta$ στ) $\beta,\zeta$.

\ex Οι $\varepsilon_1,\varepsilon_2$ είναι παράλληλες και $\alpha=70^\circ$ στο σχήμα της θεωρίας. Βρες και τις επτά άλλες γωνίες.

\ex Αν $\varepsilon_1\parallel\varepsilon_2$ και $\delta=118^\circ$, βρες τις $\beta,\zeta,\theta,\varepsilon,\alpha$ και εξήγησε με ποια ιδιότητα βρίσκεις την καθεμία.

\ex Δύο εντός εναλλάξ γωνίες παραλλήλων είναι $3x-10^\circ$ και $x+50^\circ$. Βρες το $x$ και τις γωνίες.

\ex Δύο εντός και επί τα αυτά γωνίες παραλλήλων είναι $x$ και $4x$. Βρες τες.

\ex Σωστό ή λάθος; α) Οι εντός εναλλάξ γωνίες είναι πάντα ίσες. β) Αν οι ευθείες είναι παράλληλες, οι εκτός εναλλάξ γωνίες είναι ίσες. γ) Οι εντός και επί τα αυτά γωνίες παραλλήλων είναι ίσες. δ) Δύο κατακορυφήν γωνίες είναι πάντα ίσες, είτε οι ευθείες είναι παράλληλες είτε όχι. ε) Αν οι εντός εναλλάξ γωνίες είναι ίσες, οι ευθείες είναι παράλληλες.

\ex Σχεδίασε δύο παράλληλες ευθείες που απέχουν $3$ cm και μια τέμνουσα που σχηματίζει $55^\circ$ με τη μία. Υπολόγισε όλες τις γωνίες και μέτρησε τες με τον μοιρογνωμόνιο για να ελέγξεις.

\zoni{Το χρησιμοποιώ}{κήπος, δρόμος, σχολείο}

\ex \textbf{Το μονοπάτι.} Στον κήπο της θείας Μαρίας οι δύο σειρές ντομάτας είναι παράλληλες. Μία από τις γωνίες ανάμεσα στο μονοπάτι και μια σειρά είναι $62^\circ$. Πόσες από τις οκτώ γωνίες είναι $62^\circ$ και πόσες είναι $118^\circ$;

\ex \textbf{Ο ποδηλατόδρομος.} Δύο παράλληλες λωρίδες ποδηλατόδρομου τέμνονται από μια διάβαση πεζών που σχηματίζει με την πρώτη λωρίδα μια γωνία $48^\circ$, προς το εσωτερικό του δρόμου. Πόση είναι η εντός εναλλάξ γωνία της στη δεύτερη λωρίδα; Και πόση η εντός και επί τα αυτά;

\ex \textbf{Το κιγκλίδωμα.} Σε ένα κιγκλίδωμα δύο οριζόντιες ράβδοι είναι παράλληλες και μια λοξή ράβδος τις διασχίζει. Μια γωνία της λοξής με την πάνω ράβδο, προς το εσωτερικό, είναι $125^\circ$. Πόση είναι η γωνία της με την κάτω ράβδο, επί τα αυτά μέρη και εντός;

\ex \textbf{Ο φράχτης του Βασίλη.} Ο Βασίλης λέει ότι τα δύο δοκάρια του φράχτη του σχολείου είναι παράλληλα. Μια λοξή σανίδα σχηματίζει με αυτά εντός εναλλάξ γωνίες $74^\circ$ και $76^\circ$. Έχει δίκιο;

\ex \textbf{Ο πίνακας της Φανής.} Η Φανή ελέγχει αν οι δύο απέναντι πλευρές ενός πίνακα είναι παράλληλες. Με μια λοξή ράβδο μετρά εντός και επί τα αυτά γωνίες $98^\circ$ και $82^\circ$. Είναι παράλληλες; Δικαιολόγησε.

\zoni{Το σκέφτομαι}{γιατί συμβαίνει}

\ex Εξήγησε γιατί, αν οι εντός εναλλάξ γωνίες δύο παραλλήλων είναι ίσες, τότε οι εντός και επί τα αυτά είναι παραπληρωματικές.

\ex Στο σχήμα, $\mathrm{ΒΕ}$ και $\mathrm{ΕΓ}$ είναι δύο τμήματα που ξεκινούν από το $\mathrm{Ε}$, με $\mathrm{Β}$ στην $\varepsilon_1$ και $\mathrm{Γ}$ στην $\varepsilon_2$, και $\varepsilon_1\parallel\varepsilon_2$. Η $\mathrm{ΕΒ}$ σχηματίζει με την $\varepsilon_1$ γωνία $40^\circ$ και η $\mathrm{ΕΓ}$ με την $\varepsilon_2$ γωνία $30^\circ$. Βρες τη γωνία $\widehat{\mathrm{ΒΕΓ}}$. (Βοήθεια: φέρε από το $\mathrm{Ε}$ παράλληλη προς τις δύο ευθείες.)

\begin{center}
\begin{tikzpicture}[scale=0.9]
\draw[thick] (-0.8,1.4)--(5.2,1.4) node[right] {$\varepsilon_1$};
\draw[thick] (-0.8,-1.4)--(5.2,-1.4) node[right] {$\varepsilon_2$};
\draw[thick] (3.4,1.4)--(1.73,0)--(4.15,-1.4);
\fill (3.4,1.4) circle (1.6pt) node[above] {$\mathrm{Β}$};
\fill (1.73,0) circle (1.6pt) node[left] {$\mathrm{Ε}$};
\fill (4.15,-1.4) circle (1.6pt) node[below] {$\mathrm{Γ}$};
\node at (2.5,1.15) {\small $40^\circ$};
\node at (3.2,-1.15) {\small $30^\circ$};
\end{tikzpicture}
\end{center}

\ex Στο παραλληλόγραμμο που σχηματίζουν οι παράλληλες $\varepsilon_1\parallel\varepsilon_2$ και $\varepsilon_3\parallel\varepsilon_4$, μια γωνία είναι $70^\circ$. Βρες τις τέσσερις γωνίες του τετραπλεύρου που περικλείεται.

\ex Μια οξεία και μια αμβλεία γωνία που σχηματίζονται από παράλληλες και τέμνουσα είναι εντός και επί τα αυτά, και η αμβλεία είναι τριπλάσια της οξείας. Βρες και τις δύο.

\zoni{Για το σπίτι}{ρίξε μια ματιά}

\ex Βρες στο σπίτι ή στον δρόμο δύο παράλληλες ευθείες (σιδηροτροχιά, κάγκελα, ράφια) που τέμνονται από μια τρίτη. Μέτρησε με τον μοιρογνωμόνιο δύο εντός εναλλάξ γωνίες. Είναι ίσες;

\ex Σχεδίασε με χάρακα και γνώμονα δύο παράλληλες ευθείες και μια τέμνουσα. Βάψε με το ίδιο χρώμα τις γωνίες που είναι ίσες και με άλλο χρώμα τις παραπληρωματικές τους.
